Category theory can be roughly defined as the study of mathematical
structures. If we look at the structure of the objects of study in several
fields of mathematics, we quickly notice that, in some intuitive sense, many of
these structures are quite similar. For example, most of the time the objects of
study are sets equipped with some additional elements, like binary operations
satisfying certain properties, multiplication by elements of a given field of
scalars, a collection of subsets that are considered open, etc. In addition, we
typically study the class of the functions between those sets that ``play well''
with the additional elements, generally called homomorphisms, continuous
functions, etc., to such an extent that, whenever a function is considered, the
first question to ask is, does this function belong to the class of functions of
interest? Going deeper, there are constructions on those objects that are,
intuitively speaking, ``equivalent'' to another construction in another kind of
objects, even if the structures of the objects involved are nothing
alike. Category theory attempts to formalize and thus clarify these intuitions.

Category theory was introduced by mathematicians Samuel Eilenberg and Saunders
Mac Lane in a paper published in 1942, where they introduced several of the core
concepts of the field as part of their studies in homological
algebra.\cite[p. 29]{maclane} Category theory has also been referred to as
``abstract nonsense'' or ``general abstract nonsense'', a term coined by
mathematician Norman Steenrod to affectionately refer to the kinds of very
abstract constructions and methods developed in category theory and related
fields of mathematics.\cite[preface, p. x]{riehl}

Nevertheless, the study of abstract nonsense is useful for a number of
reasons. First of all, it provides a common vocabulary for the different fields
of mathematics, which is useful for learning as having the same name for the
``same'' concept allows to better develop an intuition of the concept. Category
theory is usually not covered by undergraduate curricula, but because the other
subjects are studied with a consistent vocabulary, students grasp the idea of
concepts that are actually category-theoretic.

Second, it serves as a tool for mathematical thought. Because similar concepts
in different fields often have similar properties, it is natural for someone
with a good understanding of structures to relate to ask oneself if something
that happens in another field of mathematics, or another category, also happens
in the one being considered.

Third, because category theory is not about intuitions but about formalisms, it
provides mechanisms to transfer mathematical knowledge from one context to
another, as well as general, abstract theorems that, when appropriately made
concrete, allow proving a great variety of results. One example of this is the
Yoneda lemma, which can be used to prove that every group can be viewed as a
subgroup of a permutation proof but also that every row operation on matrices
can be defined by left multiplication by a square matrix.

Fourth, the methods of reasoning developed in category theory can sometimes be
applied to other fields

Finally, when a new field or a new kind of mathematical object is discovered, or
invented, it is often possible to use category theory to guide the study of this
new kind of object and to ``automatically'' derive properties for it. An example
of this is explored in appendix \ref{ch:programming}. Most of computer science
has been developed after category theory, so there are many concepts that have
been developed by applying general concepts of category theory to a suitable
category. To be fair, a lot of computer science theory is independent from
category theory, either because it was developed before category theory was
widely known by mathematicians or because it was not developed by mathematicians
in the first place, and a lot of it is simply unrelated to
categories. Nevertheless, as soon as the usefulness of category theory to write
and analyze computer programs was realized, a lot of its concepts, both simple
and advanced, have been imported and successfully applied to create programs
with less bugs and that are easier to maintain, as well as to design programming
languages that support these concepts and formal reasoning tools that use
category theory to support the task of verifying correctness of programs in
critical systems.

The basic object of study of category theory is a category, which consists of a
collection of objects, a collection of morphisms, operations domain ($\dom$) and
codomain ($\cod$) that take morphisms to objects, an identity operation ($1$)
that takes objects to morphisms, and a composition operation ($\circ$) that
takes pairs of morphisms $(g,f)$ with $\dom g=\cod f$ to morphisms. Furthermore,
these operations have to satisfy some properties. First, the domain and codomain
of an object's identity are both equal to the object. Second, the domain of a
composition is that of the second operand, and its codomain is that of the first
operand. Third, the composition is associative. And finally, the composition of
a morphism with the identity of its domain, as well as that of the identity of
its codomain with the morphism, are both equal to the original morphism. This
definition is independent of set theory as it only relies on predicate logic,
although we generally interpret it over set theory for convenience.

Categories represent particular mathematical concepts. Here objects might be
sets with some structure and morphisms might be functions between those sets
with certain properties, so we might have categories for e.g. groups,
topological spaces, fields, vector spaces, Banach spaces, or even for all
sets. However, objects need not be sets and morphisms need not be functions, so
we could have a category where objects are the natural numbers and morphisms
from one number to another (that is, morphisms with the first number as the
domain and the second one as the codomain) are matrices (over some ring) with so
many columns and rows, respectively. Then identities of objects would be the
corresponding identity matrices and composition would be the matrix product.

By abstracting away from the way objects and morphisms are represented in a
given category, and looking only at the relationships between those, we can
define constructions in the same way for several categories, thus formalizing
the notion of these constructions being intuitively equivalent. For example, the
product of a family a objects is defined as an object along with a family of
morphisms, called the projections, going from the product object to each object
in the family, such that for any other family of morphisms from any object to
the family of objects, there is a morphism from the domain object to the product
that, when composed with the projections, yields this family of morphisms. This
defines the Cartesian product of sets, the product topological space, the
product of groups, of ring modules, and the infimum of a subset of a
partially-ordered set.

If we ``invert'' the morphisms in a category by swapping domain and codomain and
inverting the direction of composition, the products of this category are called
coproducts of the original one, and they represent the disjoint union of sets,
the disjoint union topological space, the direct sum of groups or ring modules,
and the supremum of a subset of a partially-ordered set. It turns out that many
useful concepts can be derived from others this way.

These definitions can often be described and reasoned about visually with
commutative diagrams. If we represent morphisms as arrows between its domain and
its codomain, we can draw diagrams that look like graphs where vertices are
objects of the category and edges are morphisms, and we say that a diagram
commutes if any two paths between two vertices yield the same morphism when all
the morphisms across each path are composed together. If the morphisms are
functions, we can fix an element of a commutative diagram and ``propagate'' its
value following arrows to get a desired equality, a method known as diagram
chasing.

In order to relate concepts from different categories we use functors, which are
functions between the objects and morphisms of the two categories that respect
the operations. A nice aspect of category theory is that it is reflexive, in the
sense that it can be used to analyze itself, so we can define a category whose
objects are smaller categories and whose morphisms are functors and we
automatically get concepts like products of categories.

We can even define categories where the objects are functors between two given
categories, and then the morphisms are natural transformations. Natural
transformations are means to relate the images of two functors in a way that
plays well with the functors. Whenever some operation in a category is
considered ``natural'', or the ``natural'' way to do it, natural transformations
are usually involved. The go-to example is the standard isomorphism between a
finite-dimensional vector space and its double dual. This concept is so
important that, according to Mac Lane (\cite[p. 18]{maclane}), `` `category' has
been defined in order to define `functor' and `functor' has been defined in
order to define `natural transformation' ''.

Functors allow us to define the concept of limit, which generalizes that of
product as well as other concepts such as that of equalizer, which itself
extends the algebraic concept of kernel of a homomorphism. Similarly we have the
concept of colimits, which extends concepts like coproducts and quotient sets
and spaces.

Limits and colimits are themselves a special case of universal arrows in a
category of functors. Roughly speaking, universal arrows formalize the general
idea of an entity that satisfies some property and such that any other entity
satisfying the property is related to that one entity in a unique way. The most
common instance of this rather abstract description is that of a free object in
a construct, such as a the free group generated by a set or the free $R$-module
of a given dimension.

Universal arrows are defined in terms of a functor, and when every object in the
codomain of the functor admits a universal arrow, we can build an adjunction, a
bidirectional ``bridge'' between two categories that satisfies some coherence
properties.

In turn, adjunctions can be used to reason about monads, abstract structures
made up of an endofunctor---a functor from a category to itself---and two
natural transformations that allow to move, via morphisms, from an object to its
image under the functor, and to ``coalesce'' repeated applications of an
endofunctor over an object into only one application. Monads have two main
applications. The first one, most relevant to mathematics in general, is that
they can be used to describe a lot of common algebraic categories from other
categories, typically the category of sets. Since the category of sets is
well-understood, monads can be used to reason about these other categories. The
second application, more relevant to computer science, is that it allows
reasoning about the endofunctor itself. Endofunctors in programming are often
used to provide a ``context'' for computations, and monads can be used to thread
computations over that context by leveraging the two natural transformations.

This work is guided by two main goals. The first goal is to understand and
explain the main concepts of category theory and relate them to as many other
fields of mathematics as reasonably possible, illustrating the unifying
character of category theory. The second goal is to show how several concepts of
category theory can be applied in practice to computer science, and in
particular to programming language design and formal verification. This has been
the main motivation for covering monads, as they are incredibly useful to model
computations and interactivity in particular, even though the topic of monads is
worth studying by itself for the reasons mentioned above. On the other hand,
since the concrete explanation of how monads are used in programming is more
computer-scientific and not as mathematical in nature, this topic has been
included as an appendix.

% TODO Informal definition of category (mention that they don't have to be sets
% but they commonly are) and then rough ideas of how it works. Finally look at
% functional programming, hint at how the application was posterior to category
% theory and how polymorphic types are examples of non-constructs

% \cite{joyofcats}
% \cite{maclane}
%%% Local Variables:
%%% mode: latex
%%% TeX-master: "main"
%%% End:
